\section{Introduction}\label{sec:intro}

A Jacobi matrix whose off-diagonal entries grow fast enough is not determined by its
action on finitely supported vectors. If the entries $a_n$ grow regularly like $n^\alpha$ with
$\alpha>1$, without a strong alternation between even and odd entries, and the diagonal is small,
every solution of the eigenvalue equation is square-summable (\Cref{prop:zero} gives the precise
range): the matrix is in Weyl's \emph{limit-circle} case, the associated Hamburger moment problem is
indeterminate, and the self-adjoint realisations form a circle $\{H_t\}_{t\in[0,\pi)}$
\cite{Akhiezer1965,Simon1998}. Each $H_t$ has discrete spectrum, and the spectra of different
extensions interlace. Such matrices arise wherever an unbounded operator is written in a basis in
which it is tridiagonal: birth--death processes with polynomial rates and their Stieltjes problems
\cite{BergValent1994,Valent1999}, the higher-order squeezing operators $a^k+a^{\dagger k}$ of quantum
optics for $k\ge3$ \cite{FBL2025}, and, as a motivating example from representation theory, the
triple product of the generators of three $\mathfrak{su}(1,1)$ representations, which in a recoupling
basis is a zero-diagonal Jacobi matrix whose links grow cubically (we do not use this example below).
In each case the physically meaningful object is a particular extension, and the question is what
its spectrum is.

\paragraph{What was known.} The general theory controls the growth of the spectrum but not
its position. The eigenvalues of every $H_t$ are the zeros of an entire function of the
Nevanlinna matrix, and all these functions have the same order and type
\cite{BergPedersen1994}. For power-like links the order is $1/\alpha$
\cite{BergSzwarc2017,Romanov2017}, and in one concrete case, $a_n=n^p$, Bochkov determined the
type \cite{Bochkov2021}. General results bound the counting function of the eigenvalues
only up to universal multiplicative constants, among them the recent growth estimates of Pruckner,
Reiffenstein and Woracek \cite{PRW2023,Reiffenstein2025}; see the survey \cite{ReiffensteinWoracek2025},
whose Part~V discusses exactly this gap. Eigenvalue laws with an explicit constant were
known only for explicitly solvable chains: the quartic birth--death rates of Berg and
Valent, whose Krein and Friedrichs spectra are known in closed form
\cite{BergValent1994,Valent2006,BergChenIsmail}, and the leading term for some cubic rates
\cite{GLV2005,GLRV2006}. Valent's conjecture on the indicator of the Nevanlinna matrix
\cite{Valent1999,BergSzwarc2017} predicts the leading constant for polynomial rates, but
says nothing about the next term, which is the one that distinguishes the extensions.

\paragraph{What this paper proves.} We prove a Bohr--Sommerfeld law with an explicit offset for
every self-adjoint extension of a zero-diagonal Jacobi matrix with links
\begin{equation}\label{eq:links-intro}
  a_n=c\,n^\alpha\Bigl(1+\frac{\kappa+\kappa'(-1)^n}{n}+O(n^{-1-\delta})\Bigr),\qquad\alpha>1,
\end{equation}
which is in the limit-circle case exactly when $2|\kappa'|<\alpha-1$. For balanced links
($\kappa'=0$) the law comes with a rate, and for shifted powers with a curvature term we also
compute the second-order term. For dimerised links ($\kappa'\ne0$) the law is proved with an $o(1)$
remainder only; its rate and second-order term are open (\Cref{sec:open}). Put
$N_\lambda:=(\lambda/2c)^{1/\alpha}$, the leading-order turning scale, defined by
$2c\,N_\lambda^\alpha=\lambda$ (the site where $2a_n=\lambda$ differs from it by a bounded amount), and
\[
  \Phi(\lambda):=N_\lambda I_\alpha-\frac{\kappa\pi}{2\alpha},\qquad
  I_\alpha=\tfrac12\B\bigl(\tfrac{\alpha-1}{2\alpha},\tfrac12\bigr),\qquad
  J_\alpha=\tfrac1{2\alpha}\B\bigl(\tfrac12+\tfrac1{2\alpha},\tfrac12\bigr).
\]
The extensions are parametrised by an angle $t$ in a normalisation, fixed by the solutions
at $\lambda=0$, in which $t$ enters the law additively (\Cref{thm:count}); it is related
to Simon's and Yafaev's parametrisations by explicit formulas (\Cref{rem:kac}).

\begin{theorem}{Main theorem}{main}
Let $J$ have links \eqref{eq:links-intro} with $2|\kappa'|<\alpha-1$ and $b\equiv0$, and let
$\lambda_j(t)$ be the eigenvalues of $H_t$, labelled by $\psi(\lambda_j(t))=j\pi+t$, where $\psi$ is the lifted
phase of the regular solution at infinity (\Cref{thm:count});
for $b\equiv0$ this means $\lambda_0(0)=0$ and $\lambda_j(t)$ increases with $j$ and $t$. As $j\to+\infty$:
\begin{enumerate}
\item[(a)] if $\kappa'=0$, then $\Phi(\lambda_j(t))=j\pi+t+O(j^{-\min(1,\delta)})$ for every $t$;
\item[(b)] if $\kappa'\ne0$, put $\theta_0:=\kappa'\pi/(2(\alpha-1))$. Then
  $\Phi(\lambda_j(0))=j\pi+\theta_0+o(1)$ and $\Phi(\lambda_j(\frac\pi2))=j\pi+\frac\pi2-\theta_0+o(1)$, and
  every other extension follows the law of $H_{\pi/2}$ if $\kappa'>0$ and of $H_0$ (shifted by
  $\pi$ for $t>\frac\pi2$) if $\kappa'<0$;
\item[(c)] if $a_n=c(n+s)^\alpha(1+\kappa_2(n+s)^{-2}+O(n^{-3}))$, then
  \[
    \Phi(\lambda_j(t))=j\pi+t+\Bigl(\kappa_2+\frac{2\alpha-1}{24}\Bigr)\frac{J_\alpha}{N_{\lambda_j}}+O\bigl(N_{\lambda_j}^{-4/3}\bigr),
    \qquad\kappa=\alpha s .
  \]
\end{enumerate}
Parts (a) and (b) hold with a diagonal $|b_n|\le C(1+n)^\beta$, $\beta<\alpha-1-2|\kappa'|$: in (a) with the
labels $\psi(\lambda_j(t))=j\pi+t$ (the convention $\lambda_0(0)=0$ is then no longer available) and the
rate $O(j^{-\min(1,\delta,\alpha-1-\beta)})$, and in (b) with the offsets holding modulo $\pi$.
\end{theorem}

Part (a) is \Cref{cor:rate}, which needs no condition on the differences of the links; part (b) is
\Cref{cor:dimer} with \Cref{thm:dimerdiag}; part (c) is \Cref{thm:second}; the counting functions
follow (\Cref{cor:balanced}). Three features of the law are worth stating separately.

First, there is no quarter-period correction. Below the turning point each two-site block contributes
$\pi$ to the phase. The Airy turning point contributes its quarter-period $+\frac\pi4$, and the
limit-circle end contributes $-\frac\pi4$: the turning-point mode takes equal values on the two sites of
a block, and the solutions that define the boundary conditions were normalised to equal amplitude at
infinity, so in their coordinates the mode sits at $-\frac\pi4$ (\Cref{sec:turning}, \Cref{fig:budget}).
The statement that does not depend on any coordinate on the circle of extensions is the law for $H_0$,
the extension defined by the solution $P$ of the equation at $\lambda=0$ with $P_0=1$, $P_1=0$:
rescaling $P$ does not change the boundary condition, and for balanced links
$\Phi(\lambda_j(0))=j\pi+o(1)$. The same holds for $H_{\pi/2}$, defined by the second solution $R$, with
$j\pi+\frac\pi2$. Every other extension needs a coordinate, and the normalisation of $P$ and $R$ to
equal amplitude at infinity (by the constant $\vp$ of \Cref{prop:zero}) is the one in which the
band-edge mode sits at exactly $-\frac\pi4$ for every $\lambda$, which is why $t$ then enters additively.
In Simon's coordinate $\tau_{\rm S}$, related by $\tan t=\tan\tau_{\rm S}/\vp$, the same law is not additive. So the absence
of a quarter-shift is a statement relative to this calibration, which is fixed once by the data at
$\lambda=0$ and not adjusted to any eigenvalue.

Second, the offset $\kappa/(2\alpha)$ is set by the $1/n$ term of the links, and the extension enters only
through $t/\pi$.

Third, an alternating $1/n$ term changes the far field qualitatively. The two solutions at $\lambda=0$
acquire block envelopes of different sizes, the equation becomes a Bessel problem of order
$\nu=\frac12-\kappa'/(\alpha-1)$, and the extensions split into the one defined by the solution of
smaller block-envelope size ($H_0$ if $\kappa'>0$, $H_{\pi/2}$ if $\kappa'<0$) and all the others. The
coordinate-free form of the dimerised law is that the offsets of $H_0$ and $H_{\pi/2}$, the extensions
defined by $P$ and by $R$, move from their balanced values $0$ and $\frac\pi2$ by $+\theta_0$ and
$-\theta_0$, $\theta_0=\kappa'\pi/(2(\alpha-1))$; this is visible already in the leading behaviour of the
eigenvalues. The second-order constant also has an interpretation: $\kappa_2J_\alpha$ is the classical
correction from the curvature of the links, and $(2\alpha-1)/24=(\alpha^2-(\alpha-1)^2)/24$ is the
Weyl--Moyal term of the symbol $2A(X)\cos p$ (\Cref{rem:moyal}).

\paragraph{What is new and what is not.} The leading constant is the one predicted by
Valent's conjecture under regular distribution of the zeros, and the completely regular
growth and indicator of the Nevanlinna matrix that follow from the law (\Cref{cor:type})
can also be obtained from Bochkov's type theorem together with the Berg--Szwarc reduction
and a Tauberian argument. The limit-circle criterion $2|\kappa'|<\alpha-1$ is due to \'Swiderski
and Trojan \cite{SwiderskiTrojan2022}. The factorisation of the transfer matrix into shears
is the classical factorisation of the Nevanlinna matrix and the transfer matrix of Kac's
canonical system \cite{Kac1999}. What we believe is new is the $O(1)$ term and everything
beyond it: the offsets and their dependence on the extension, the Bessel far field for
dimerised links and its consequences (including the Krein and Friedrichs offsets
$(\dsum+\esum)/(2p^2)\mp\nu/2$ for all polynomial birth--death rates, \Cref{cor:bd}), the rates with
their sharpness, the second-order term, and the exact critical indices of \Cref{thm:valent}.
These indices bear on two questions of Berg and Szwarc \cite{BergSzwarc2026}. Their open problem, whether both of their classes $\mathbf A$ and $\mathbf B$ of N-extremal solutions
are nonempty, is answered for $\mathbf B$ by explicit measures (\Cref{thm:valent}); for $\mathbf A$ it follows from the
case $r=2$ of a theorem of Borichev and Sodin, by a construction that is not explicit
(\Cref{cor:AB}). We also locate the determinacy threshold of \cite[Ex.~4.1]{BergSzwarc2026} at
$c=-2$. We found no eigenvalue law with an explicit offset for any class of limit-circle Jacobi
matrices beyond the explicitly solvable chains.

\paragraph{Method.} The proof works with one monotone function, the lifted phase $\psi$ of
the regular solution at infinity, whose level sets are the spectra of all extensions
(\Cref{sec:setting}). Three devices compute it. An exact identity turns each two-site block
of the shear product into a rotation by $\varphi_m=2\arcsin(\lambda/2a_{2m+1})$ in a slowly
varying frame, so that an adiabatic normal form sums the frame changes (\Cref{sec:blocks});
a compactness argument handles a window of $O(N_\lambda^{1/3})$ blocks at the turning point
(\Cref{sec:turning}); and, for dimerised links, a Bessel comparison handles the far field
(\Cref{sec:dimer}). For rates and the second-order term we prove a discrete analogue of
Olver's error bounds for Airy approximations (\Cref{thm:olver}), an exact inequality whose only weights
are Olver's envelopes. The same bound answers a question of Fischer, Burgarth and Lonigro: their
description of essentially singular limits of Jacobi operators holds for growth exponents $\alpha>1$,
not only $\alpha>\frac43$ (\Cref{thm:fbl}). Finally, the Bessel boundary values determine the residues
of the Nevanlinna functions of a birth--death problem at their poles. Combined with a weighted
determinacy criterion, the case $r=2$ of a theorem of Borichev and Sodin specialised to N-extremal
supports (\Cref{prop:weighted}), this gives the critical indices $(1\mp\nu)/2$ of all N-extremal
solutions of Valent's indeterminate problems. Berg and Szwarc introduced these indices and proved
$a(\mu_F)\le\frac12\le a(\mu_\sigma)$ for Stieltjes problems, exhibiting members of the classes at the
value $\frac12$; no index had been computed (\Cref{thm:valent}).

\paragraph{Organisation.} \Cref{sec:setting} fixes notation and reduces the spectral problem
to the phase $\psi$. \Cref{sec:blocks,sec:turning} prove the block identity, the normal form
and the turning-point analysis. \Cref{sec:balanced} states the first-order law for balanced
links and its consequences, and \Cref{sec:dimer} treats dimerised links and birth--death
processes. \Cref{sec:olver} proves the discrete Olver bound, which \Cref{sec:rates,sec:second}
use for rates and the second-order term. \Cref{sec:apps} applies the results to squeezing
operators and to singular limits, \Cref{sec:moments} to the critical indices of indeterminate
moment problems, \Cref{sec:numerics} compares every headline number with
an independent computation, and \Cref{sec:open} lists open problems. \Cref{app:routine}
contains the technical estimates and auxiliary proofs, so that the main text can follow the ideas,
and \Cref{app:repro} the computational cross-checks and the code and data availability.
\Cref{sec:history} says how the work began and how AI tools were used, and
\Cref{app:provenance} how the results were reached.
